\section*{अध्याय \ref{ch:g9:metals-acids-corrosion} --- धातुएँ: अम्लों के साथ अभिक्रियाएँ और संक्षारण}

\begin{solution}{exo:g9:metals-acids-corrosion:1}
हाइड्रोजन: नली के मुँह पर \omterm{def:g4:what-a-fire-needs:burning}{जलती} खपच्ची तीखी `पॉप' की आवाज़ देती है।
\end{solution}

\begin{solution}{exo:g9:metals-acids-corrosion:2}
लोहा, ज़िंक और ऐलुमिनियम अभिक्रिया करते हैं; ताँबा और सोना नहीं।
\end{solution}

\begin{solution}{exo:g9:metals-acids-corrosion:3}
ऑक्सीजन और पानी, दोनों।
\end{solution}

\begin{solution}{exo:g9:metals-acids-corrosion:4}
ऐसा \omterm{def:g9:ions:ion}{आयन} जो अभिक्रिया के दौरान मौजूद रहता है पर उसमें भाग नहीं लेता।
हाइड्रोक्लोरिक अम्ल के साथ: क्लोराइड \omterm{def:g9:ions:ion}{आयन} \ce{Cl-}।
\end{solution}

\begin{solution}{exo:g9:metals-acids-corrosion:5}
\ce{Mg(s) + 2H+(aq) -> Mg^{2+}(aq) + H2(g)}; हर ओर आवेश $+2$।
\end{solution}

\begin{solution}{exo:g9:metals-acids-corrosion:6}
सोडियम हाइड्रॉक्साइड मिलाओ: \ce{Fe(OH)2} का हरा \omterm{def:g9:ions:precipitate}{अवक्षेप}
आयरन(II) \omterm{def:g9:ions:ion}{आयन} दिखाता है।
\end{solution}

\begin{solution}{exo:g9:metals-acids-corrosion:7}
उबालने से पानी में घुली \omterm{def:g4:air-a-mixture-of-gases:air}{वायु} निकल जाती है, और तेल \omterm{def:g4:air-a-mixture-of-gases:air}{वायु} को फिर से \omterm{def:g2:mixing-and-dissolving:dissolve}{घुलने} से रोकता
है: कील के पास पानी है पर ऑक्सीजन नहीं।
\end{solution}

\begin{solution}{exo:g9:metals-acids-corrosion:8}
कार पर उछलकर पड़ने वाले पानी में घुला नमक \omterm{def:g9:metals-acids-corrosion:corrosion}{जंग} लगने को तेज़ कर देता है।
\end{solution}

\begin{solution}{exo:g9:metals-acids-corrosion:9}
बनी हुई ऑक्साइड की परत धातु से चिपक जाती है और उसे \omterm{def:g4:air-a-mixture-of-gases:air}{वायु} से बंद कर देती है:
आक्रमण सतह पर ही रुक जाता है।
\end{solution}

\begin{solution}{exo:g9:metals-acids-corrosion:10}
उसे पेंट करो, तेल या ग्रीस लगाओ, यशदलेपित करो (या उसे स्टेनलेस स्टील का बनाओ)।
\end{solution}

\begin{solution}{exo:g9:metals-acids-corrosion:11}
\ce{2Al(s) + 6H+(aq) -> 2Al^{3+}(aq) + 3H2(g)}: हर ओर 2 Al और 6 H,
हर ओर आवेश $+6$। 200 Al \omterm{def:g7:atoms-and-molecules:atom}{परमाणुओं} के लिए: $200 \times \frac{3}{2} = 300$
हाइड्रोजन \omterm{def:g7:atoms-and-molecules:molecule}{अणु}।
\end{solution}

\begin{solution}{exo:g9:metals-acids-corrosion:12}
खरोंच के आस-पास का ज़िंक स्टील की जगह संक्षारित होता है: जब तक पास में ज़िंक
बचा है, स्टील बचा रहता है।
\end{solution}

\begin{solution}{pb:g9:metals-acids-corrosion:1}
\textbf{1.} \ce{Zn(s) + 2H+(aq) -> Zn^{2+}(aq) + H2(g)}।

\textbf{2.} हाइड्रोजन: \omterm{def:g4:what-a-fire-needs:burning}{जलती} खपच्ची के साथ तीखी `पॉप' की आवाज़।

\textbf{3.} सोडियम हाइड्रॉक्साइड \ce{Zn(OH)2} का सफ़ेद
\omterm{def:g9:ions:precipitate}{अवक्षेप} देता है।

\textbf{4.} क्लोराइड \omterm{def:g9:ions:ion}{आयन} \ce{Cl-}।

\textbf{5.} $125.6 - 113.5 = \qty{12.1}{g}$।

\textbf{6.} जब तक ज़िंक अम्ल के साथ अभिक्रिया करता है, हाइड्रोजन निकलती है;
ज़िंक न बचने पर बुलबुले रुक जाते हैं।

\textbf{7.} स्टील का लोहा भी अभिक्रिया करता, हाइड्रोजन और
\ce{Fe^{2+}} \omterm{def:g9:ions:ion}{आयन} देता; तब सोडियम हाइड्रॉक्साइड हरा \omterm{def:g9:ions:precipitate}{अवक्षेप} देता।

\textbf{8.} स्टील की जगह ज़िंक पर आक्रमण होता है, और वह पानी और \omterm{def:g4:air-a-mixture-of-gases:air}{वायु} को स्टील
से दूर रखता है।

\textbf{9.} $2 \times 10.0 \times 10.0 = \qty{200}{cm^2}$।

\textbf{10.} $12.1 \div 7.13 \approx \qty{1.70}{cm^3}$।

\textbf{11.} $1.70 \div 200 = \qty{0.0085}{cm}$।

\textbf{12.} $0.0085 \times 10\,000 = \qty{85}{\micro m}$: ज़िंक की
परत लगभग \qty{85}{\micro m} मोटी थी।
\end{solution}
